% Options for packages loaded elsewhere
\PassOptionsToPackage{unicode}{hyperref}
\PassOptionsToPackage{hyphens}{url}
\PassOptionsToPackage{dvipsnames,svgnames,x11names}{xcolor}
%
\documentclass[
]{article}
\usepackage{amsmath,amssymb}
\usepackage{iftex}
\ifPDFTeX
  \usepackage[T1]{fontenc}
  \usepackage[utf8]{inputenc}
  \usepackage{textcomp} % provide euro and other symbols
\else % if luatex or xetex
  \usepackage{unicode-math} % this also loads fontspec
  \defaultfontfeatures{Scale=MatchLowercase}
  \defaultfontfeatures[\rmfamily]{Ligatures=TeX,Scale=1}
\fi
\usepackage{lmodern}
\ifPDFTeX\else
  % xetex/luatex font selection
\fi
% Use upquote if available, for straight quotes in verbatim environments
\IfFileExists{upquote.sty}{\usepackage{upquote}}{}
\IfFileExists{microtype.sty}{% use microtype if available
  \usepackage[]{microtype}
  \UseMicrotypeSet[protrusion]{basicmath} % disable protrusion for tt fonts
}{}
\makeatletter
\@ifundefined{KOMAClassName}{% if non-KOMA class
  \IfFileExists{parskip.sty}{%
    \usepackage{parskip}
  }{% else
    \setlength{\parindent}{0pt}
    \setlength{\parskip}{6pt plus 2pt minus 1pt}}
}{% if KOMA class
  \KOMAoptions{parskip=half}}
\makeatother
\usepackage{xcolor}
\usepackage[margin=25mm]{geometry}
\usepackage{longtable,booktabs,array}
\usepackage{calc} % for calculating minipage widths
% Correct order of tables after \paragraph or \subparagraph
\usepackage{etoolbox}
\makeatletter
\patchcmd\longtable{\par}{\if@noskipsec\mbox{}\fi\par}{}{}
\makeatother
% Allow footnotes in longtable head/foot
\IfFileExists{footnotehyper.sty}{\usepackage{footnotehyper}}{\usepackage{footnote}}
\makesavenoteenv{longtable}
\setlength{\emergencystretch}{3em} % prevent overfull lines
\providecommand{\tightlist}{%
  \setlength{\itemsep}{0pt}\setlength{\parskip}{0pt}}
\setcounter{secnumdepth}{-\maxdimen} % remove section numbering
\ifLuaTeX
  \usepackage{selnolig}  % disable illegal ligatures
\fi
\usepackage{bookmark}
\IfFileExists{xurl.sty}{\usepackage{xurl}}{} % add URL line breaks if available
\urlstyle{same}
\hypersetup{
  pdftitle={Matched Shot Distributions: Scientific Design and Verification},
  colorlinks=true,
  linkcolor={Maroon},
  filecolor={Maroon},
  citecolor={Blue},
  urlcolor={Blue},
  pdfcreator={LaTeX via pandoc}}

\title{Matched Shot Distributions: Scientific Design and Verification}
\author{}
\date{}

\begin{document}
\maketitle

\section{Matched Shot Distributions: Scientific Design and
Verification}\label{matched-shot-distributions-scientific-design-and-verification}

\subsection{Purpose and Estimands}\label{purpose-and-estimands}

This extension compares straight, draw, and fade patterns with a common
\textbf{physical zero-error landing point}, followed by prescribed
delivery noise. No simulated endpoint or trajectory is rotated to aim at
the target. Nominal face direction, path, and speed are inputs solved
before sampling. Matching two endpoint coordinates does not match the
entire trajectory, launch state, peak height, descent angle, or energy.
Curved trajectories remain different by design.

The study estimates changes in landing spread, bias, supported-shot
probability, and conditional course scoring under explicit equipment and
delivery assumptions. It does not identify a universal preferred shape,
a measured human variability distribution, or the performance of a
particular commercial club.

This is an editable research reference, separate from the governed
engineering manual. Its \href{matched_delivery_reference.tex}{Standalone
LaTeX Version} is regenerated from this research Markdown with Pandoc,
not from the engineering manual. The \href{ASTRA_REVIEW.md}{Astra Design
Review}, \href{../README.md}{Phase-One Reference}, issue records in
\url{issues/}, executable tests, and immutable result receipts form its
evidence chain. Implementation and results remain under development
until their acceptance evidence is recorded.

\subsection{Experimental Controls}\label{experimental-controls}

For each club, use static loft and lie from the cited specification,
fixed attack angle, and explicit assumed head properties. Pitch the
complete club by forward lean of 0°, 5°, or 10°. Then solve nominal face
and speed, with path constrained by face-to-path difference: 0° for
straight, −1.5° for draw, +1.5° for fade; the stronger curves use ±3°.
The target line remains the fixed +X direction throughout.

Two target protocols distinguish the questions:

\begin{enumerate}
\def\labelenumi{\arabic{enumi}.}
\tightlist
\item
  \textbf{Within-Lean Matching:} Straight, draw, and fade land at the
  same target for a given club and lean. The target is the zero-error
  straight shot at that club\textquotesingle s reference speed and lean.
  Distances may differ across lean settings.
\item
  \textbf{Across-Lean Matching:} Attempt the same physical target
  distance as the zero-lean straight reference. Report all infeasible
  conditions with their fixed bounds and reason. Reachability is not
  guaranteed by an optimizer\textquotesingle s termination.
\end{enumerate}

The across-lean Monte Carlo control uses the shaft-only and
combined-2-mph profiles at both twist SDs and both curve strengths (72
planned cells). It holds target distance constant where matching
succeeds. It is a matched-delivery lean comparison: matching also
changes nominal speed, face and path, and a fixed absolute speed SD
gives a different coefficient of variation at different nominal speeds.
It does not identify a pure isolated lean effect or establish results
for the seven other noise profiles.

Speed search bounds are fixed at 75--125\% of the reference speed;
nominal face bounds remain local to the intended shape. The accepted raw
landing residual must be at most 0.05 m, checked independently after
optimization. The frozen experiment plan records actual bounds and
solved inputs. Infeasible cases remain in the matrix; neither attack
angle nor static loft is silently changed to populate them.

\subsection{Equipment and Provenance}\label{equipment-and-provenance}

\begin{longtable}[]{@{}lrrrrr@{}}
\toprule\noalign{}
Club & Static Loft & Lie at Zero Lean & Reference Speed & Attack Angle &
Head Mass \\
\midrule\noalign{}
\endhead
\bottomrule\noalign{}
\endlastfoot
Driver & 10° & 58.5° & 45 m/s & −0.9° & 0.200 kg \\
7-Iron & 34° & 63° & 36 m/s & −4° & 0.272 kg \\
Pitching Wedge & 46° & 64° & 32 m/s & −5° & 0.300 kg \\
\end{longtable}

The \href{https://www.titleist.com/golf-clubs/drivers/gt2}{Titleist GT2
specification} provides the selected driver loft/lie. The
\href{https://www.titleist.com/golf-clubs/irons/720-t100}{historical
T100 specification} provides the iron/wedge lofts and lies. These
references establish geometry only. Speed, attack, mass, CG, inertia,
usable face dimensions, restitution, friction, and curvature are
research assumptions. The hybrid presets are \textbf{not}
reconstructions or performance predictions for GT2 or T100 clubs.

\begin{longtable}[]{@{}lrlll@{}}
\toprule\noalign{}
Club & CG Depth & Principal Inertias: Normal, Heel, Up & Face Half-Width
/ Half-Height & Bulge / Roll Radius \\
\midrule\noalign{}
\endhead
\bottomrule\noalign{}
\endlastfoot
Driver & 35 mm & 0.00025, 0.00040, 0.00045 kg·m² & 50 / 27 mm & 300 /
300 mm \\
7-Iron & 10 mm & 0.00012, 0.00020, 0.00025 kg·m² & 35 / 23 mm & Flat /
Flat \\
Pitching Wedge & 8 mm & 0.00010, 0.00020, 0.00022 kg·m² & 35 / 25 mm &
Flat / Flat \\
\end{longtable}

Restitution is 0.83 and friction coefficient 0.4 in the primary
experiment. Inertia and CG sensitivity must accompany off-center
conclusions. \href{https://doi.org/10.1016/j.proeng.2016.06.329}{McNally
et al. (2016)} establishes the importance of CG depth, inertia, and face
geometry for off-center trajectory sensitivity; it does not validate
these assumed parameter values. Manufacturer loft alone is insufficient
to identify gear effect.

\subsection{Frames and Rigid Delivery
Geometry}\label{frames-and-rigid-delivery-geometry}

The native frame is right-handed: +X toward target, +Y toward the
right-handed golfer/left, +Z upward. Reported lateral coordinates and
face/path angles are positive right. Contact horizontal offset is
positive toward the toe (−face-heel axis); vertical offset is positive
toward the top of the face.

Starting with the static face normal and zero-lean shaft axis, apply a
rigid rotation about +Y for forward lean, then a negative world-Z
rotation for positive-right nominal face direction. Apply the stochastic
shaft rotation using Rodrigues\textquotesingle{} formula about that
delivered shaft axis. Positive shaft rotation closes the right-handed
face. Zero twist gives delivered loft = static loft − forward lean.
These rotations preserve the shaft--face angle and orthonormal face
frame.

With static loft L, lie β, lean ℓ, nominal right-positive face f, and
shaft twist θ, the geometric construction is

\[\mathbf n_0=(\cos L,0,\sin L)^T,\qquad
\mathbf s_0=(0,\cos\beta,\sin\beta)^T,\]

\[Q=R_z(-f)R_y(\ell),\quad \mathbf s=Q\mathbf s_0,\quad
\mathbf n=R_{\mathbf s}(\theta)Q\mathbf n_0.\]

Rodrigues\textquotesingle{} rotation is
\(R_{\mathbf s}(\theta)=I\cos\theta+[\mathbf s]_\times\sin\theta+(1-\cos\theta)\mathbf s\mathbf s^T\).
The reported face and loft are \(\operatorname{atan2}(-n_y,n_x)\) and
\(\operatorname{atan2}(n_z,\sqrt{n_x^2+n_y^2})\). Internal trigonometry
uses radians; public angles use degrees. This explicitly explains why
shaft rotation can change both face and loft and why forward lean
changes that coupling. It does not establish an observed
player\textquotesingle s closure trajectory.

At zero shaft twist, the first-order sensitivities follow directly from
\(d\mathbf n/d\theta=\mathbf s\times\mathbf n\):

\[\frac{df}{d\theta}=-\frac{\sin\beta\cos L}{\cos(L-\ell)},\qquad
\frac{d\lambda}{d\theta}=-\cos\beta,\qquad
\frac{d\lambda}{df}=\cot\beta\,\frac{\cos(L-\ell)}{\cos L}.\]

Here λ is delivered loft. Ratios are unchanged when all angle increments
use degrees. Consequently, fixed shaft-twist SD produces approximately
the same first-order loft SD across lean settings, while face-angle SD
changes. This does not mean carry must increase on every closed-face
miss: launch height, spin, ball speed and flight response also change.
Finite-angle results and actual landing slopes are computed by the full
model; these derivatives describe local geometric sensitivity only. The
\href{validation/GEOMETRY_DERIVATIVES.md}{Independent Geometry and Carry
Derivative Checks} provide the analytical derivation, finite-difference
evidence and club-specific limitations.

Unlike phase one\textquotesingle s constant-delivered-loft sensitivity,
forward lean here changes both face loft and shaft orientation. A 10°
static driver at 10° lean can have zero delivered loft and a nonairborne
nominal launch; that is a physical-domain limitation to report, not a
missing value to interpolate.

\subsection{Two-Body Contact Law}\label{two-body-contact-law}

The previous off-center path combined a scalar effective mass with
additive heuristic spin and omitted the corresponding club rotational
impulse. It is unsuitable for claiming mechanically consistent strike
dispersion. The new research path uses a full club inertia tensor and
contact lever, without additive gear spin.

Let c be club-minus-ball relative contact velocity, n the unit outward
local face normal, r\_c and r\_b the CG-to-contact levers, and {[}r{]}×
the cross-product matrix:

\[\mathbf c=\mathbf v_c+\boldsymbol\omega_c\times\mathbf r_c-\mathbf v_b-\boldsymbol\omega_b\times\mathbf r_b,\qquad \mathbf r_b=-R\mathbf n.\]

The delivery adapter assumes zero initial club angular velocity and a
stationary, nonspinning ball. Shaft twist is an orientation error, not a
closure angular speed. Contact-generated club rotation is retained after
impact. For impulse J applied to the ball:

\[K=\left(\frac{1}{m_c}+\frac{1}{m_b}\right)I
-[\mathbf r_c]_\times I_c^{-1}[\mathbf r_c]_\times
-[\mathbf r_b]_\times I_b^{-1}[\mathbf r_b]_\times,\qquad
\mathbf c_{\rm after}=\mathbf c_{\rm before}-K\mathbf J.\]

K has units kg⁻¹, impulse J has units N·s, and relative velocity c has
units m/s. The ball is a solid sphere in this idealization. The full
tensor is rotated from the specified principal frame into world
coordinates before the impulse calculation.

\[\Delta\mathbf v_b=\mathbf J/m_b,\quad
\Delta\boldsymbol\omega_b=I_b^{-1}(\mathbf r_b\times\mathbf J),\quad
\Delta\mathbf v_c=-\mathbf J/m_c,\quad
\Delta\boldsymbol\omega_c=-I_c^{-1}(\mathbf r_c\times\mathbf J).\]

The primary law integrates Coulomb contact in normal-impulse
coordinates. For each positive normal increment dP, solve the tangential
maximum-dissipation disk problem using the current relative contact
velocity and the full coupled K. Compression ends at zero normal
approach velocity. Store compression work E\_c = ∫ c\_n dP; restitution
returns e\_energy² E\_c through ∫ −c\_n dP. The coefficient 0.83 is
therefore \textbf{energetic restitution}, not a prescribed normal
velocity ratio in coupled off-center contact. Resolution is explicit and
must pass refinement checks. The optional Numba implementation uses no
fast-math transformations and is checked against the Python reference.

For an orthonormal tangent basis T and current relative velocity c, the
incremental disk problem is

\[\min_{\|\Delta\mathbf J_t\|\leq\mu\,dP}
\frac12\Delta\mathbf J_t^T(T^TKT)\Delta\mathbf J_t
-\left[T^T(\mathbf c-dP\,K\mathbf n)\right]^T\Delta\mathbf J_t.\]

The total increment is \(dP\,\mathbf n+T\Delta\mathbf J_t\). Normal work
uses trapezoidal integration, with bounded bisection at the compression
and restitution endpoints. The primary resolution is 128 compression
steps; source-bound refinement to 256/512 steps and flight dt
0.02/0.01/0.005 s is recorded in the
\href{validation/independent_preflight_summary.json}{Independent
Numerical Evidence}.

A separate maximum-dissipation/Newton comparator minimizes ½JᵀKJ − cᵀJ
subject to nᵀKJ = (1+e)c\_n and the impulse friction cone. This
comparator does not claim to reproduce the full slip/grip history. An
earlier stick-first candidate passed simple centered checks but violated
passivity in Astra\textquotesingle s coupled elastic counterexample;
that failed experiment motivated the corrected comparator and
incremental primary law. It is not used in the final primary pathway.

All four body velocity states receive equal-and-opposite linear and
angular impulse updates. Validate work recovery, cone bounds, energy
passivity, linear/angular momentum, frame covariance, and numerical
refinement; do not replace failed numerical solutions with heuristic
launches. These rigid contact laws are not measured
tangential-compliance models. Conservation and convergence establish
implementation consistency, not empirical golf-ball/club accuracy. Shaft
flexibility, face compliance, speed-dependent restitution, local
friction, grooves, and turf interaction remain unqualified. The external
measurement obligation remains open as
\href{https://github.com/D-sorganization/UpstreamDrift/issues/11984}{\#11984}
and in the
\href{../../../development/planning/DV-11984.md}{Deferred-Validation
Catalog}.

For curved faces, a local parabolic surface provides a declared geometry
approximation. Surface displacement along the center normal is
−½(h²/R\_b + v²/R\_r), and its normal is proportional to n +
(h/R\_b)t\_toe + (v/R\_r)t\_up. A zero radius explicitly selects a flat
direction; it is not evaluated as a division by zero.

\subsection{Deterministic Head-Assumption Sensitivity}
A separate deterministic sensitivity check holds face, path and shaft twist at
zero, retains preset speed and attack angle, and does not rematch targets.
It examines all three clubs and lean settings, with center strikes and
0.5\,in toe, heel, high and low offsets. Of 300 variant evaluations, 268 are
supported and 32 are unsupported. The offsets are fixed test positions,
not the stochastic study's 0.25\,in per-axis strike standard deviation.

\begin{center}
\begin{tabular}{lrrrr}
\toprule
Club & CG Depth & Inertia & Face Radii & Contact Law\\
\midrule
Driver & 5.14 & 7.10 & 22.63 & 11.33\\
7-Iron & 0.54 & 1.76 & --- & 1.22\\
Pitching Wedge & 0.18 & 1.40 & --- & 0.33\\
\bottomrule
\end{tabular}
\end{center}
Values are maximum landing displacements in metres among comparable supported
baseline/variant pairs. CG depth, principal inertias together, and driver
bulge/roll radii together are varied by $\pm20\%$. These deliberately chosen
ranges are not measured confidence intervals. The last column instead changes
contact law: the primary energetic restitution and comparator Newton velocity
coefficient both use the number 0.83 but have different definitions. Maxima
can occur at different deliveries; they are not additive or dispersion SDs.
The driver's largest face-radius response occurs at $5^\circ$ lean with a low
strike and radius scale 0.8, illustrating sensitivity at low launch.

The \href{validation/independent_head_sensitivity_summary.json}{saved
summary} and \href{validation/independent_head_sensitivity.json}{full
execution receipt} preserve these cases and their unsupported outcomes.
This model dependence is distinct from finite-sample Monte Carlo uncertainty;
small statistical intervals do not establish empirical impact-model accuracy.

\subsection{Stochastic Delivery and
Covariance}\label{stochastic-delivery-and-covariance}

The primitive error vector is ordered as shaft rotation (degrees), path
(degrees), speed (m/s), toe/heel contact (m), high/low contact (m),
forward lean (degrees), and attack angle (degrees). A multivariate
Gaussian is a transparent local sensitivity assumption, not an
established distribution for all human motion. Without player data there
is no justified empirical tail or covariance fit.

Default comparisons include shaft-rotation SD 1° and 2°, path SD 0.5°,
speed SD 1 and 2 mph (exact conversion 0.44704 m/s per mph), and strike
SD 0.25 inch = 0.00635 m on \textbf{each} contact axis. The radial
strike SD is not 0.25 inch. Baseline and progressively combined noise
profiles separate contributions before their interactions are
considered.

Users can edit all SDs and the complete symmetric correlation matrix.
Require a unit diagonal, finite entries, and positive semidefiniteness.
Zero-SD dimensions are exactly constant. Draw one seeded standardized
sample matrix and reuse corresponding shot indices across patterns.
Preserve the requested covariance and report the realized input
covariance. Actual face angle is derived nonlinearly from shaft rotation
and lean; a shaft/path correlation must not be labelled an exact
face/path correlation. Report realized face/path covariance explicitly.

Do not clip, retry, or resample off-face, nonairborne, or otherwise
unsupported deliveries. Preserve their input rows and failure
categories. Supported-endpoint statistics and density are explicitly
conditional when any shots are unsupported; report the valid fraction
and interval alongside them. Full-population strokes gained is
unavailable if an unsupported outcome has no defensible terminal lie.
Conditional paired scoring uses the intersection of supported indices
across all patterns, with the reduced population stated.

\subsection{Density, Uncertainty, and Course
Scoring}\label{density-uncertainty-and-course-scoring}

Academic figures use common axes, units, and density normalization
across patterns. Show points, binned probability density, marginal
distributions, and empirical 50\%, 80\%, and 95\% highest-density
regions. Report bin resolution and achieved probability mass; tied bins
can exceed the requested threshold. Do not manufacture a two-dimensional
Gaussian width for a nearly one-dimensional input/output manifold.
Export publication PNG, SVG, and PDF; inspect rendered output for
clipping, overlaps, and readable annotations.

Report lateral and longitudinal SD, full covariance, mean bias, radial
RMSE and quantiles, target containment, and paired differences with
sampling uncertainty. The same approximate historical scoring baseline
as phase one applies: driver tee/fairway/rough; iron/wedge
approach/green/rough. Landing is a hypothetical stopping point. Roll,
hazards, and wedge stopping are not inferred. Sampling intervals exclude
model, equipment, and course uncertainty. Course-dimension and
off-center head-parameter sensitivities remain separate from Monte Carlo
error.

\subsection{Pilot, Final Matrix, and
Acceptance}\label{pilot-final-matrix-and-acceptance}

First run small deterministic and Monte Carlo pilots. The final
10,000-shot-per-pattern matrix begins only after software contracts,
physical identities, matching residuals, numerical refinement,
population accounting, plot review, and an independent Astra review
pass. Save unavailable cells as explicit records rather than omitting
them.

\begin{longtable}[]{@{}lll@{}}
\toprule\noalign{}
Requirement & Executable Evidence & Acceptance \\
\midrule\noalign{}
\endhead
\bottomrule\noalign{}
\endlastfoot
Static Loft and Lean & \texttt{test\_club\_geometry.py},
\texttt{test\_matched\_config.py} & Rotation identities, physical
ranges, sourced static values \\
Correlated Noise & \texttt{test\_noise.py} & PSD rejection,
reproducibility, zero variance, empirical covariance \\
Raw Nominal Matching & \texttt{test\_nominal\_matching.py} and Real
Pilot Receipts & Physical raw X/Y residual ≤0.05 m; fixed bounds; no
aiming rotation \\
Off-Center Contact & Independent Rigid Contact Tests & Momentum, passive
energy, restitution, frame covariance, central limit, gear direction \\
Unsupported Samples & Matched Runner Tests & Every input retained; no
resampling; explicit conditional estimands \\
Density Figures & Density/Plot Tests and Rendered Review &
Normalization, ties, degeneracy, common scale, readable exports \\
Final Results & Source-Bound Manifests and Independent Recalculation &
Exact reproducibility, complete matrix accounting, no stale
provenance \\
\end{longtable}

Final reports must link exact source revisions, native binary identity,
input specifications, seeds, solver bounds, numerical tolerances, result
hashes, pilot evidence, and review dispositions. Historical phase-one
outputs remain labelled with their original model and controls.

\end{document}
